\section{Introduction}
A scene query asks for an identity and its spatial extent. Text may name an
object, an image may show one of its parts, and a point may identify one visible
surface. A useful three-dimensional response must connect that evidence to the
requested object while rejecting distractors and recovering members beyond the
observed region. Strong teacher features provide valuable correspondences, but
a high match score alone does not specify the extent of a correct answer.

Two decisions are easily conflated. First, an observation may or may not refer
to the requested identity. Second, even a correctly matched observation may
support only part of that identity's extent. Treating its exterior as background
can erase unseen object parts; accepting an unrelated region can propagate
false positives through otherwise accurate geometry. The problem is therefore
not only how to aggregate evidence, but which evidence is entitled to support
or exclude each part of the response.

We study \emph{query-conditioned evidence fusion}: native observations are
admitted according to their applicability to the query, assigned spatial
support and exclusion authority, and combined through a fixed scene carrier.
Rejected or unavailable evidence remains unknown. Text applicability uses
candidate-relative native matching; reference tracks, categorical competition
and explicit points retain their respective conditioning semantics. This
response-level interface uses existing teacher representations without learning
a common latent feature field.

Evidence formation must also preserve the intended output capability.
Continuous semantic responses and complete selected instances have different
optimization objectives and readouts. We retain both, using geometric response
fusion for semantic evidence and complete-membership mechanisms for actual
selection. This avoids interpreting silhouette membership as a calibrated
semantic probability, or assuming that a better score field automatically
improves an independent selected set. The shared evidence contract has typed
implementations; it does not require an identical solver for every query type.

Our implementation, recorded as v10 for reproducibility, covers LERF semantic
segmentation and selected-set projection, ScanNet open-vocabulary segmentation,
NVOS reference segmentation and multimodal UQIS. Controlled component studies
connect identity admission, admitted-response use, local image correspondence
and complete membership to measured improvements. Further matched experiments
separate changes in evidence construction, reconstruction, response decisions
and output interpretation. These controls also expose cases where better source
fit fails to improve target answers. External comparisons retain their declared
teacher, geometry and input-access differences.

Our contributions are threefold. First, we formulate a query-conditioned
evidence contract that separates identity applicability, spatial support,
exclusion and missing information. Second, we implement this contract in a
scene-query method that combines native evidence through geometry while
preserving both semantic response and complete selection. Third, we evaluate
five task outputs with matched component controls, protocol checks and complete
positive and negative feedback. Optional response-replacement and NVOS
extensions are evaluated separately from the established system, so their
measured effects can be attributed to the operations they actually change.
